\section*{الفصل \ref{ch:b3:radicals-carbenes} --- الجذور والكربينات وإعادات الترتيب}

\begin{solution}{exo:b3:radicals-carbenes:1}
للبروبان ست ذرات هيدروجين أولية وذرتان ثانويتان. الكلورة: $6 \times 1 : 2 \times 3.9 = 6 : 7.8$،
أي 43\,\% من 1-كلوروبروبان و57\,\% من 2-كلوروبروبان. البرومة: $6 : 164$، أي 4\,\% 
و96\,\%.
\end{solution}

\begin{solution}{exo:b3:radicals-carbenes:2}
$\ce{CH2}$: ثلاثي؛ $\ce{CCl2}$: أحادي (فالأزواج الحرة للكلور ترفع المدار p الفارغ).
ويُضاف ثنائي كلوروكربين إضافة نوعية فراغيًا: \textit{trans}-1،1-ثنائي كلورو-2،3-ثنائي ميثيل سيكلوبروبان.
\end{solution}

\begin{solution}{exo:b3:radicals-carbenes:3}
\textit{cis}-1،2-ثنائي إيثيل سيكلوبروبان: يُضاف \omterm{def:b3:radicals-carbenes:carbenoid}{الكربينويد} في خطوة واحدة إلى وجه واحد.
\end{solution}

\begin{solution}{exo:b3:radicals-carbenes:4}
السيكلوهكسانون: $\varepsilon$-كابرولاكتون (أوكسيبان-2-ون). الأسيتوفينون: أسيتات الفينيل،
$\ce{CH3CO2C6H5}$؛ فالفينيل يهاجر قبل الميثيل.
\end{solution}

\begin{solution}{exo:b3:radicals-carbenes:5}
تسع ذرات هيدروجين أولية وذرة ثالثية واحدة. الكلورة: $9 : 5.1$، أي 64\,\% من 1-كلورو-2-ميثيل بروبان
و36\,\% من 2-كلورو-2-ميثيل بروبان. البرومة: $9 : 1600$، أي 0.6\,\% و99.4\,\%.
\end{solution}

\begin{solution}{exo:b3:radicals-carbenes:6}
$k_c = 3.0 \times k_H[\mathrm{Bu_3SnH}] = 3.0 \times \num{2.4e6} \times 0.020 = \qty{1.4e5}{s^{-1}}$.
\end{solution}

\begin{solution}{exo:b3:radicals-carbenes:7}
الهجوم على C5: 5-exo-trig، مفضل؛ وعلى C6: 6-endo-trig، مفضل بحسب القواعد لكنه
أبطأ هنا. 4-exo-tet: مفضل. 5-endo-trig: غير مفضل. 5-exo-dig: مفضل.
3-exo-dig: غير مفضل. 5-endo-dig: مفضل.
\end{solution}

\begin{solution}{exo:b3:radicals-carbenes:8}
يعطي فقد الماء كاتيونًا ثالثيًا بنزيليًا؛ وعلى الكربون المجاور يمكن أن تنتقل مجموعة فينيل
أو مجموعة ميثيل، فيهاجر الفينيل، الأعلى قابلية. ثم يفقد الكاتيون
المستقر بالأكسجين بروتونًا: 3،3-ثنائي فينيل بوتان-2-ون.
\end{solution}

\begin{solution}{exo:b3:radicals-carbenes:9}
كورتيوس: $\ce{C6H5CON3 -> C6H5NCO + N2}$. ومع الماء، يفقد حمض الكرباميك $\ce{CO2}$: الأنيلين
(ويتكوّن قليل من ثنائي فينيل اليوريا حين يلتقي الأنيلين بمزيد من \omterm{def:b3:radicals-carbenes:curtius}{الإيزوسيانات}). ومع الإيثانول:
\textit{N}-فينيل كاربامات الإيثيل، $\ce{C6H5NHCO2C2H5}$.
\end{solution}

\begin{solution}{exo:b3:radicals-carbenes:10}
الكلور: الأولي $421.7 - 432 = \qty{-10}{kJ/mol}$، والثالثي $405.4 - 432 = \qty{-27}{kJ/mol}$؛ والبروم: $+56$ 
و$+39$. فانتزاعا الكلور كلاهما ناشران للحرارة، بحالتين انتقاليتين مبكرتين
لا تحسان إلا قليلًا بالفرق البالغ \qty{16}{kJ/mol}؛ وانتزاعا البروم كلاهما ماصان للحرارة،
بحالتين انتقاليتين متأخرتين تختلف طاقتا تنشيطهما بمعظمه:
$\eu^{16000/(8.314 \times 298)} \approx 600$، وهي رتبة النسبة المرصودة.
\end{solution}

\begin{solution}{exo:b3:radicals-carbenes:11}
الأوكسيم \textit{E} (المجموعة OH في وضع anti من C2، الحامل لمجموعة الميثيل): يهاجر C2، ويدخل النيتروجين
بجواره: 7-ميثيل أزيبان-2-ون. الأوكسيم \textit{Z}: يهاجر C6: 3-ميثيل أزيبان-2-ون. ففي
\omterm{def:b3:radicals-carbenes:beckmann}{إعادة ترتيب بكمان} تحسم هندسةُ الأوكسيم الأمرَ؛ وفي \omterm{def:b3:radicals-carbenes:baeyer}{أكسدة باير--فيليغر}
تحسمه \omterm{def:b3:radicals-carbenes:migratory}{القابلية للهجرة}، فينتقل الكربون الثانوي:
7-ميثيل أوكسيبان-2-ون.
\end{solution}

\begin{solution}{exo:b3:radicals-carbenes:12}
النسبة المتحلّقة $k_c/(k_c + k_H c)$: 0.91 و0.49 و0.09. ومن أجل 95\,\%: $k_Hc = k_c(1/0.95 - 1)$، ومنه
$c = 0.0526 \times \num{2.3e5}/\num{2.4e6} = \qty{5.0e-3}{mol/L}$. ويُحافَظ عليه بإضافة الهيدريد ببطء
بمضخة محقنة، أو باستعمال كمية حفزية من هاليد القصدير مع عامل مختزل
يجدد الهيدريد.
\end{solution}

\begin{solution}{pb:b3:radicals-carbenes:1}
\textbf{1.} $\ce{C6H10O + NH2OH -> C6H11NO + H2O}$.

\textbf{2.} يلزم الحمض لإزالة مجموعة الهيدروكسي من وسيط الإضافة في صورة
ماء، لكن الحمض الزائد يبرتن الهيدروكسيل أمين، فلا يعود يُضاف.

\textbf{3.} لا: فذرتا الكربون المرتبطتان بكربون الرابطة C=N متكافئتان بحكم
تناظر الحلقة.

\textbf{4.} في الأوكسيم \textit{E} تتجه المجموعة OH بعيدًا عن الكربون الحامل للميثيل C2 (في وضع anti
منه)، وفي الأوكسيم \textit{Z} نحوه.

\textbf{5.} للانقلاب عند نيتروجين الأوكسيم الحامل لأكسجين كهرسلبي حاجز
مرتفع.

\textbf{6.} يبرتن مجموعة الهيدروكسي (أو يسلفنها)، فيجعلها مجموعة مغادرة
جيدة، أي الماء.

\textbf{7.} تنتقل الرابطة C--C في الحلقة الواقعة في وضع anti من المجموعة المغادرة من الكربون إلى النيتروجين،
الذي يُدرج بذلك في الحلقة: فتصير الذرات الست سبعًا.

\textbf{8.} أيون نتريليوم حلقي.

\textbf{9.} $\ce{C6H11NO -> C6H11NO}$: تماكب، من أوكسيم إلى لاكتام.

\textbf{10.} الكابرولاكتام، أزيبان-2-ون، الأميد الحلقي السباعي لحمض
6-أمينوهكسانويك؛ وبلمرته بفتح الحلقة تعطي النايلون-6.

\textbf{11.} \qty{98.0}{g/mol} و\qty{113.0}{g/mol}.

\textbf{12.} $\qty{1.0e6}{g}/\qty{98.0}{g/mol} = \qty{10204}{mol}$؛ $\times 0.98 \times 0.98 = \qty{9800}{mol}$؛ $\times \qty{113.0}{g/mol} = \qty{1.11e6}{g}$، أي \qty{1.11}{t}.

\textbf{13.} الأوكسيم: $\num{10204} \times 0.98 = \qty{1.00e4}{mol}$؛ و$\ce{H2SO4}$: $\qty{1.30e4}{mol}$، تعطي العدد نفسه من مولات
$\ce{(NH4)2SO4}$: $\qty{1.30e4}{mol} \times \qty{132.1}{g/mol} = \qty{1.7}{t}$، أي أكثر من الكابرولاكتام.

\textbf{14.} $\ce{C6H10O + RCO3H -> C6H10O2 + RCOOH}$.

\textbf{15.} ناتج الإضافة الرباعي الأوجه للحمض البيروكسي على كربون الكربونيل (وسيط
كريغي)، وهو هيمي كيتال بيروكسي.

\textbf{16.} تنتقل مجموعة $\ce{CH2}$ من الحلقة (والمجموعتان متكافئتان) إلى أقرب أكسجين
بيروكسيدي؛ ويغادر الكربوكسيلات $\ce{RCOO-}$.

\textbf{17.} $\varepsilon$-كابرولاكتون، أوكسيبان-2-ون؛ وبلمرته بفتح الحلقة تعطي
بولي($\varepsilon$-كابرولاكتون)، وهو بوليستر قابل للتحلل الحيوي.

\textbf{18.} ناتجه الجانبي الوحيد هو الماء، بينما يترك الحمض البيروكسي مولًا من
الحمض الكربوكسيلي لكل مول من اللاكتون ويكون هو نفسه خطرًا في التخزين.

\textbf{19.} يهاجر الكربون الثانوي C2: 7-ميثيل أوكسيبان-2-ون.

\textbf{20.} نعم: يحتفظ الكربون المهاجر بتشكيله.

\textbf{21.} 7-ميثيل أزيبان-2-ون (C2 في وضع anti، فيهاجر).

\textbf{22.} 3-ميثيل أزيبان-2-ون (يهاجر C6).

\textbf{23.} بكمان: هندسة الأوكسيم (تنتقل المجموعة الواقعة في وضع anti).
باير--فيليغر: \omterm{def:b3:radicals-carbenes:migratory}{القابلية للهجرة} (ينتقل الكربون الأكثر استبدالًا).

\textbf{24.} ذرتا الكربون $\alpha$ فيه متكافئتان: أوكسيم واحد، ولاكتام واحد، ولاكتون واحد.

\textbf{25.} يعطي طن واحد من السيكلوهكسانون نحو \qty{1.11}{t} من الكابرولاكتام.
\end{solution}
